\documentclass[10pt,a4paper]{article}
\usepackage[utf8]{inputenc}
\usepackage[T1]{fontenc}
\usepackage[french]{babel}
\usepackage{amsmath}
\usepackage{amsfonts}
\usepackage{amssymb}
\usepackage{titlesec}
\usepackage{minted}
\author{Alexandre Moine}
\title{Rapport de projet}

\newcounter{question}

\titleformat{\subsection}
{\normalfont\normalsize\bfseries}{Question \stepcounter{question}\thequestion}{1em}{}

\begin{document}
\maketitle

\section{Introduction}
% aussi présent pour le numbering

\section{Entrées} % TODO: trouver comment rajouter "appetizers"

\subsection{}
\begin{minted}{ocaml}
  predicate no_3_consecutive_zeros_sub (a:array int) (l:int) =
  forall i. 0 <= i < l-2 -> not (a[i] = a[i+1] = a[i+2] = 0)
\end{minted}
La borne supérieure du prédicat est $l-2$ car en supposant $0 < l <= a.length$, on veut que $i+2 < a.length$.

\subsection{}
\begin{minted}{ocaml}
  predicate no_3_consecutive_zeros (a:array int) =
  no_3_consecutive_zeros_sub a a.length
\end{minted}
En effet, il n'y a pas 3 zéros consécutifs dans $a$ si et seulement s'il n'a pas 3 zéros consécutifs entre les indices $0$ et la taille de $a$.

\subsection{}
\begin{minted}{ocaml}
for i = 0 to a.length - 3 do
  invariant { no_3_consecutive_zeros_sub a (i+2) }
    if a[i] = 0 && a[i+1] = 0 && a[i+2] = 0 then raise TripleFound
\end{minted}
Les bornes de la boucles sont claires: on s'assure que les indices utiles au test sont valides.\\
Pour l'invariant, on voit que la fonction lève une exception lorsqu'elle trouve 3 zéros consécutifs entre $i$ et $i+2$. À chaque fin de tour, on a donc qu'il n'y a pas 3 zéros consécutifs entre l'indice $0$ et $i+2$.
\subsection{}
\begin{minted}{ocaml}
for i = 2 to a.length - 1 do
  invariant { last1 = a[i-1] /\ last2 = a[i-2] }
  invariant { no_3_consecutive_zeros_sub a i }
    let v = a[i] in
    if v = 0 && last1 = 0 && last2 = 0 then raise TripleFound;
    last2 <- last1; last1 <- v;
\end{minted}
Les bornes de la boucles sont claires: on s'assure que les indices utiles au test sont valides.\\
Pour l'invariant, on voit que la fonction lève une exception lorsqu'elle trouve 3 zéros consécutifs entre $i-2$ et $i$. À chaque fin de tour, on a donc qu'il n'y a pas 3 zéros consécutifs entre l'indice $0$ et $i$.

\subsection{}
\begin{minted}{ocaml}
let rec ghost function num_occ (e:int) (f:int -> int) (i j :int) : int
  variant { j-i }
  = if i>=j then 0 else
    if f (j-1) = e then 1 + num_occ e f i (j-1) else num_occ e f i (j-1)
\end{minted}
On \og parcoure\fg{} la fonction de $j$ à $i$, le variant est donc $i-j$. De plus, sans précondition sur $j$, on donne un test large d'arrêt: $ i >= j$.

\subsection{}

\begin{minted}{ocaml}
let count_number_of (e:int) (a:array int) : int
  ensures { result = num_occ e (fun i -> a[i]) 0 a.length}
  = ...
\end{minted}
On veut s'assurer que la fonction compte bien le nombre d'éléments égaux à $e$ dans tout $a$.

\subsection{}
\begin{minted}{ocaml}
  let ref n = 0 in
  for i=0 to a.length - 1 do
    invariant { n = num_occ e (fun i -> a[i]) 0 i }
    if a[i] = e then n <- 1+n;
  done; n
\end{minted}
L'invariant est trivial ici: la référence $n$ compte bien le nombre d'éléments égaux à $e$ entre $0$ et $i$.

\subsection{}
\begin{minted}{ocaml}
predicate identical_sub_arrays (a:array int) (o1 o2 l:int)
= forall k:int. 0 <= k < l -> a[o1+k] = a[o2+k]
\end{minted}
Deux sous-tableaux sont égaux s'ils le sont à chaque indices.
\subsection{}
\begin{minted}{ocaml}
predicate valid_sub_array (a:array int) (o l:int) =
  0 <= o /\ o + l <= a.length

let check_identical_sub_arrays (a:array int) (o1 o2 l:int) : bool
  requires { valid_sub_array a o1 l /\ valid_sub_array a o2 l  }
  ensures { result = True <-> identical_sub_arrays a o1 o2 l }
= try
    for k=0 to l-1 do
      invariant { identical_sub_arrays a o1 o2 k }
      if a[o1+k] <> a[o2+k] then raise DiffFound
  done;
  ...
\end{minted}

On commence par se donner un prédicat $valid\_sub\_array a o l$ spécifiant que le sous-tableau $a[o... (o+l-1)]$ est valide. On requiert ensuite que les sous-tableau donné en argument soient valides.\\
L'invariant suit encore une fois directement la spécification. Au rang $i$ de la boucle, on est sûr que les deux sous-tableaux $a[o_1...(o_1+i-1)]$ et $a[o_1...(o_2+i-1)]$ sont égaux.

\section{Takuzu}
\subsection{}
\subsection{}
\subsection{}
\subsection{}
\subsection{}
\subsection{}
\subsection{}
\subsection{}
\subsection{}
\subsection{}
\subsection{}
\subsection{}
\subsection{}
\subsection{}

\section{Conclusion}
\subsection{}

\end{document}
%%% Local Variables:
%%% mode: latex
%%% TeX-master: t
%%% End:
